\documentclass[11pt,a4paper]{article}

\usepackage[utf8]{inputenc}
\usepackage[T1]{fontenc}
\usepackage[spanish]{babel}
\usepackage{lmodern}
\usepackage{geometry}
\usepackage{hyperref}
\usepackage{enumitem}
\usepackage{longtable}
\usepackage{array}
\usepackage{xcolor}
\usepackage{listings}
\usepackage{amsmath}

\geometry{margin=2.5cm}
\hypersetup{
    colorlinks=true,
    linkcolor=blue,
    urlcolor=blue,
    pdftitle={Ataque compuesto sobre wallet local en librustzcash},
    pdfauthor={Codex}
}

\lstdefinestyle{plain}{
    basicstyle=\ttfamily\small,
    breaklines=true,
    frame=single,
    columns=fullflexible
}

\title{Analisis tecnico del ataque compuesto sobre wallet local}
\author{Documento interno de auditoria tecnica}
\date{\today}

\begin{document}

\maketitle

\tableofcontents
\newpage

\section{Resumen ejecutivo}

Este documento reemplaza el informe anterior centrado en un solo bug de PCZT.
La tesis actual es mas ambiciosa:

\begin{quote}
No estamos frente a un unico bug aislado, sino frente a una cadena de bugs
wallet-side que se pueden combinar para producir un impacto local mayor:
historial falso, recipient intent falso, estado spendable corrupto,
persistencia de metadata no revalidada, fallos tardios al gastar y
denegacion de servicio durante recovery / scan / rescan.
\end{quote}

La conclusion central es:

\paragraph{Lo que no afirmamos.}
No confirmamos robo directo de fondos on-chain ni bypass de consenso.

\paragraph{Lo que si afirmamos.}
Si confirmamos una cadena plausible y fuerte de \textbf{corrupcion de estado
de wallet local}.

El ataque compuesto mas fuerte que hoy podemos sostener tecnicamente es:

primero, el wallet acepta o procesa datos shielded no plenamente validados;
despues, persiste campos de nota y metadata auxiliar sin rebindearlos
suficientemente al commitment y al output real; mas adelante, en flujos PCZT,
sigue confiando en metadata mutable para recipient, cuenta, valor y direccion
visible; luego, tras un deep rewind, puede preservar parte de ese estado
criptografico inconsistente; y finalmente varias ramas permiten panic antes de
que una correccion por rescan o reprocesamiento limpie el estado.

Eso lleva a un impacto agregado tipo:

sent history falso, received history semantica o estructuralmente
inconsistente, memo u output index mal asociados, spendability local
incorrecta, witness o nullifier state persistiendo mas de lo debido, fallos
tardios al construir spends y DoS operacional.

\section{Objetivo del documento}

Este documento persigue cuatro metas. Primero, describir el
\textbf{ataque compuesto} que queremos demostrar con PoC. Segundo, explicar
\textbf{muy detalladamente} los bugs chicos que lo habilitan. Tercero, dejar
claro cuales partes son \textbf{hallazgos confirmados} y cuales son
\textbf{composicion / amplificacion}. Cuarto, fijar el
\textbf{plan de implementacion de PoC} en la misma forma usada en la branch
anterior: como tests en \texttt{zcash\_client\_backend},
\texttt{zcash\_client\_memory} y, para las rutas wallet-side relevantes,
tambien en \texttt{zcash\_client\_sqlite}.

\section{Estado actual de PoC}

La cadena compuesta ya no es solo un diseno teorico. En esta branch quedaron
PoC ejecutables para dos rutas ya demostradas. La primera es la
\textbf{corrupcion de intent en PCZT}: hay tests backend para Sapling y
Orchard, tests sobre \texttt{zcash\_client\_memory} y ahora tambien tests
sobre \texttt{zcash\_client\_sqlite} que muestran que la sent history puede
persistir un recipient falso o reutilizar una reclasificacion falsa de output
externo como estado wallet-interno. La segunda es el
\textbf{mismatch de contexto ZIP-212}: hay un test backend que demuestra que la
misma transaccion puede ser tratada como wallet-owned o no segun el contexto
de cadena que se pase a \texttt{decrypt\_transaction}; y esa misma
superficie wallet-visible ya quedo revalidada tambien en
\texttt{zcash\_client\_sqlite} para la ruta de sent history.

En otras palabras, hoy ya tenemos demostradas por tests la ruta
``PCZT mutable $\rightarrow$ historial falso'', la ruta
``PCZT mutable $\rightarrow$ reclasificacion de output externo como
\texttt{InternalAccount}'', la ruta
``reclasificacion falsa $\rightarrow$ resumen historico reutiliza esa
semantica interna'', la ruta
``contexto legacy $\rightarrow$ aceptacion wallet-side bajo semantica
equivocada'' y la ruta
``contexto legacy $\rightarrow$ sent history visible'' tanto para el wallet
in-memory como para \texttt{zcash\_client\_sqlite} cuando la transaccion se
interpreta como \textit{outgoing} bajo el contexto equivocado.

Ademas, esta branch ya contiene una \textbf{PoC compuesta sobre el mismo
wallet} que une ambas familias en una sola historia operativa. Primero, el
wallet acepta una transaccion Sapling legacy por falta de contexto de cadena y
la expone por \texttt{get\_sent\_outputs()}. Despues, en ese mismo wallet, una
segunda transaccion creada por el flujo
proposal $\rightarrow$ PCZT $\rightarrow$ prove $\rightarrow$ sign
$\rightarrow$ extract queda reclasificada como \texttt{InternalAccount}
aunque el output comprometido siga siendo externamente recoverable. Finalmente,
el mismo wallet queda en un estado donde la tx legacy sigue mostrando
recipient externo por \texttt{get\_sent\_outputs()}, mientras que la tx PCZT
ya no muestra ningun recipient externo por esa misma API.

El efecto compuesto mas fuerte que hoy ya esta demostrado por test es que,
cuando ambas corrupciones coexisten en el mismo wallet, \texttt{get\_tx\_history()}
deja de producir una vista global consistente y falla identificando la
transaccion legacy admitida por el path sin contexto.

Esa ya no es la ultima etapa demostrada. Hoy la branch tambien contiene una
\textbf{variante ejecutable de escalada a High}: se toma ese mismo wallet ya
corrompido por las dos familias previas, se lo deja en un estado todavia
suficientemente vivo como para seguir escaneando y luego un unico bloque
compacto malformado en la siguiente altura hace panicar el
\texttt{scan\_cached\_blocks()} normal.

En la variante actual, el campo malformado elegido es \texttt{prev\_hash} con
longitud invalida, de modo que el panic aparece durante la verificacion de
continuidad de bloques en el scanner y no por un hook artificial del test.

La parte que \textbf{todavia no} esta ejecutable en
\texttt{zcash\_client\_memory} es la amplificacion por reorg profundo, porque
ese target todavia no implementa \texttt{rewind\_to\_chain\_state}.

\subsection{Estado de integracion en Zallet}

La PoC ya no vive solamente dentro de \texttt{librustzcash}. Tambien quedo
integrada en el repo de \texttt{Zallet} en dos niveles.

\paragraph{Variante white-box.}
La primera variante es \textbf{white-box fuerte} sobre el engine de sync de
Zallet. Ejerce \texttt{initialize} y el loop real de \texttt{steady\_state}
sobre un \texttt{MockChain} y demuestra panic cuando el siguiente bloque
malformado entra por la ruta normal de sync.

\paragraph{Variante black-box.}
La segunda es una \textbf{black-box de producto} que ejecuta
\texttt{zallet start} contra un validator RPC stub minimo y demuestra que el
proceso termina de manera fallida por input malicioso proveniente del origen
normal de datos de cadena.

Hay una diferencia tecnica importante entre ambas. La variante white-box sobre
el loop de sync conserva la idea del \texttt{CompactBlock} malformado. La
variante black-box por subprocess no usa el mismo \texttt{prev\_hash}
truncado, porque en la ruta
\texttt{FetchService -> BlockCache -> CompactBlock} Zaino reconstruye ese
campo desde el bloque crudo. Por eso, en \texttt{zallet start}, la superficie
efectiva elegida fue \texttt{z\_getsubtreesbyindex}: Zallet hace
\texttt{unwrap()} sobre roots devueltos por el validator durante
\texttt{initialize}, y un root malformado produce el crash de producto.

Este punto es importante para disclosure: la composicion exacta del bug puede
variar entre capas, pero el mensaje operativo sigue siendo el mismo:
\textbf{un wallet local real puede caerse por input malicioso proveniente de su
fuente normal de sync}.

La integracion en Zallet ya subio un paso mas: no solo se confirmo un crash de
producto aislado, sino tambien \textbf{persistencia del fallo tras restart}.
Hoy ya hay una PoC donde:

\begin{enumerate}[leftmargin=2em]
    \item se inicializa una wallet real;
    \item se la hace arrancar contra el mismo validator RPC stub malicioso;
    \item \texttt{zallet start} falla;
    \item sin reparar nada en disco ni cambiar de backend, se vuelve a correr
    \texttt{zallet start};
    \item y la wallet vuelve a fallar por la misma condicion.
\end{enumerate}

Eso no prueba robo ni consenso, pero si refuerza mucho el framing de
\textit{freezing} operacional: el usuario no enfrenta un crash accidental
unico, sino un estado donde la wallet no puede progresar autonomamente mientras
persista la misma fuente de sync maliciosa.

Ademas, la integracion de Zallet ya mejoro en dos direcciones adicionales que
suben la calidad probatoria del impacto. Por un lado, ya existe una variante
white-box donde el wallet tiene balance transparente visible por
\texttt{WalletSummary} antes de entrar en el crash del loop de
\texttt{steady\_state}. Por otro lado, ya existe una variante separada donde
la tarea real de \texttt{recover\_history} paniquea al procesar bloques
historicos malformados, en vez de depender del arranque o del steady-state
loop.

Eso fortalece dos partes distintas del argumento:

\begin{enumerate}[leftmargin=2em]
    \item que el crash puede afectar una wallet que ya tiene estado economico
    visible y no solo una wallet ``vacia'';
    \item que el mismo defecto no solo rompe el sync normal, sino tambien una
    ruta de recuperacion que el usuario razonablemente intentaria usar para
    volver a un estado sano.
\end{enumerate}

Es importante describir con precision el alcance actual de estas dos mejoras:

\begin{itemize}[leftmargin=2em]
    \item el caso de \textbf{balance visible antes del crash} ya esta probado
    en una variante white-box del loop real de \texttt{steady\_state}, y ademas
    ya existe una variante black-box donde \texttt{zallet start} falla contra
    el validator malicioso despues de que el \texttt{wallet.db} fue presembrado
    con un UTXO wallet-owned reconocido;
    \item el caso de \textbf{crash en recovery historico} ya esta probado en
    una variante white-box de la tarea real de \texttt{recover\_history}, y
    ademas ya existe una variante black-box por subprocess donde
    \texttt{zallet start} termina prematuramente al entrar por una historia
    historica real servida por un validator malicioso;
    \item la variante black-box por subprocess de \texttt{zallet start} hoy
    sigue demostrando la caida de producto y la persistencia tras restart;
    de las dos propiedades adicionales, ahora ya incorpora tanto la de estado
    wallet-owned previo al crash como la de historia historica real. Lo que
    cambia entre variantes es \textbf{como} se manifiesta el fallo.
\end{itemize}

En otras palabras, la evidencia actual ya demuestra las tres cosas fuertes por
separado:

\begin{enumerate}[leftmargin=2em]
    \item crash o terminacion prematura de producto;
    \item persistencia del fallo tras restart;
    \item impacto sobre wallet con estado economico visible;
    \item impacto sobre recovery historico.
\end{enumerate}

Lo que todavia queda como mejora opcional, y no como condicion para sostener el
reporte, seria reunir todas esas propiedades dentro de una unica PoC black-box
de producto con una semantica de fallo uniforme. Hoy la variante historica por
subprocess ya existe, pero se observa como terminacion prematura de
\texttt{zallet start} porque el supervisor actual descarta el
\texttt{Result} de las tareas de sync y apaga el proceso en cuanto alguna de
ellas termina.

Ese paso de unificacion tambien quedo alcanzado. Hoy ya existe una variante
black-box de producto donde, dentro del mismo escenario general:

\begin{enumerate}[leftmargin=2em]
    \item la wallet ya tiene estado wallet-owned visible en disco;
    \item el trabajo pendiente real cae en historia historica;
    \item el validator malicioso provoca la terminacion de
    \texttt{zallet start};
    \item un segundo \texttt{zallet start} vuelve a terminar de manera
    prematura bajo la misma fuente maliciosa;
    \item e incluso una reparacion local agresiva del estado wallet no evita la
    terminacion mientras persista el mismo backend.
\end{enumerate}

Con eso, las propiedades que antes estaban demostradas por piezas quedaron ya
demostradas tambien en una PoC black-box unificada:

\begin{itemize}[leftmargin=2em]
    \item estado wallet-owned visible;
    \item historia historica real;
    \item terminacion de producto;
    \item persistencia tras restart;
    \item freezing practico frente a reinicios y frente a una reparacion local
    del estado.
\end{itemize}

La unica reserva que sigue siendo importante es semantica:
\textbf{la variante historica por subprocess se manifiesta como terminacion
prematura del proceso, no necesariamente como exit code no-cero}, porque el
supervisor actual trata la salida de una tarea ongoing como condicion de
apagado aunque el error no se propague como codigo de salida.

Pero desde el punto de vista del impacto operacional sobre el usuario, esa
reserva ya no cambia demasiado la lectura: la wallet queda en una condicion de
\textit{freezing} practico mientras siga dependiendo de la misma fuente de sync
maliciosa.

\section{Fuentes usadas dentro de esta branch}

Los dos documentos auxiliares de esta branch que guiaban la PoC anterior son
los siguientes. Uno es \path{LIBRUSTZCASH_AUDIT_FINDINGS.md}. El otro es
\path{LIBRUSTZCASH_PCZT_AUDIT.md}.

De ellos extraemos dos decisiones metodologicas. La primera es que la PoC debe
vivir como \textbf{tests ejecutables} dentro del repo. La segunda es que el
lugar correcto para empezar es \texttt{zcash\_client\_backend} y
\texttt{zcash\_client\_memory}. Eso ya dio resultado, y desde ahi la misma
familia de PoC relevantes pudo bajarse tambien a
\texttt{zcash\_client\_sqlite} para confirmar que la corrupcion no era un
artefacto del backend mock ni del store in-memory.

\section{Hipotesis principal}

La hipotesis de trabajo para la PoC compuesta es:

\begin{quote}
Un atacante con capacidad de introducir estado no plenamente validado o
metadata mutable en un wallet local puede conseguir que el wallet almacene y
reutilice una vista incorrecta del estado shielded, y que esa corrupcion se
propague a historial, recipient intent, spendability local y recovery /
rescan, con posibilidad adicional de DoS.
\end{quote}

\subsection{Capacidades del atacante}

No hace falta que el atacante tenga \emph{todas} estas capacidades juntas.
El documento agrupa varias rutas. Las capacidades relevantes son un servidor
remoto malicioso o corrupto que alimenta compact blocks, input local de
recovery o import que llama APIs wallet-side sobre transacciones no plenamente
validadas, un coordinador o tooling step malicioso dentro de un flujo PCZT, o
un momento de deep rewind o rescan sobre estado ya contaminado.

\section{Mapa del ataque compuesto}

\subsection{Vista de alto nivel}

La cadena completa puede verse asi:

\[
\begin{array}{c}
\texttt{admission}
\rightarrow
\texttt{persistence}
\rightarrow
\texttt{intent\ corruption}
\\[0.3em]
\rightarrow
\texttt{spendability\ corruption}
\rightarrow
\texttt{reorg\ amplification}
\rightarrow
\texttt{panic/doS}
\end{array}
\]

Cada fase puede ocurrir o no. No hace falta completar toda la cadena para
tener impacto relevante, pero el peor caso local aparece cuando varias fases se
encadenan.

\section{Fase 1: admision de estado no validado}

\subsection{Bug 1A: \texttt{decrypt\_and\_store\_transaction} no es un validador de consenso}

\paragraph{Hallazgo confirmado.}
\texttt{decrypt\_and\_store\_transaction} puede decriptar y persistir salidas
wallet-relevant sin primero probar que la transaccion sea consenso-valida.

\paragraph{Consecuencia.}
Entra al wallet un estado que semantica y estructuralmente se trata como si ya
hubiese pasado una frontera de validacion mas fuerte de la que realmente paso.

\subsection{Bug 1B: compact scanning puede promover outputs sin validacion local completa}

\paragraph{Hallazgo confirmado.}
El path de compact scanning puede elevar outputs shielded a un estado local
que parece wallet-owned o spendable-looking, sin verificacion local de proofs
ni signatures.

\paragraph{Consecuencia.}
Esto no implica aceptacion de consenso, pero si implica que el wallet puede
formar opiniones locales fuertes sobre ownership y spendability antes de una
validacion completa.

\subsection{Bug 1C: confusion de contexto ZIP-212}

\paragraph{Hallazgo confirmado.}
Si falta contexto de altura, \texttt{decrypt\_transaction} puede terminar
aplicando el modo de enforcement pre-ZIP-212 (\texttt{Off}) a transacciones no
minadas actuales.

\paragraph{Consecuencia.}
El wallet puede aceptar plaintext semantics pre-ZIP-212 en un contexto donde
consenso actual no las aceptaria.

\subsection{Por que esta fase es importante}

La Fase 1 define la frontera de ingreso:

\begin{quote}
el wallet empieza a trabajar con datos shielded que son decryptable o
parseable, pero todavia no estan suficientemente amarrados a la semantica de
consenso que el resto del sistema informalmente asume.
\end{quote}

\section{Fase 2: persistencia sin rebinding suficiente}

\subsection{Bug 2A: nota persistida sin recheck de commitment}

\paragraph{Hallazgo confirmado.}
Una vez que una \texttt{Received*Output} cruza a storage, los componentes de
nota se persisten sin recomputar y comparar contra el \texttt{cmu}/\texttt{cmx}
publicado.

Esto rompe la idea deseable:

\[
\texttt{stored note state} \Longleftrightarrow \texttt{published commitment}
\]

porque la equivalencia deja de revalidarse en la frontera de persistencia.

\subsection{Bug 2B: nullifier y tree position caller-supplied}

\paragraph{Hallazgo confirmado.}
Nullifier y commitment-tree-position pueden aceptarse y persistirse como
metadata provista por el caller, en vez de rederivarse o cross-checkearse de
forma canonica.

\paragraph{Consecuencia.}
El wallet puede construir una vision local de spentness o witness-availability
basada en metadata no suficientemente reatada al note real.

\subsection{Bug 2C: memo y output index no reatados al mismo output cifrado}

\paragraph{Hallazgo confirmado.}
Memo y output/action index pueden persistirse como piezas independientes sin
reafirmar que siguen apuntando al mismo output ciphertext-bearing.

\subsection{Bug 2D: upserts sticky}

\paragraph{Hallazgo confirmado.}
Ciertas asociaciones wallet-side en SQLite pueden quedar ``pegadas'' bajo
semantica de upsert. En particular, la confirmacion reciente sobre
\texttt{zcash\_client\_sqlite} mostro una causa concreta de divergencia:
\texttt{sent\_notes} retenia \texttt{to\_address} y \texttt{to\_account\_id}
previos aunque una actualizacion posterior reclasificara el output.

Ese detalle importaba directamente para las PoC de reclasificacion falsa de
output externo como \texttt{InternalAccount}, porque hacia que SQLite siguiera
mostrando recipient externo en sent history y tx history incluso cuando el
estado ya habia sido reescrito como wallet-interno. Corregido ese punto, las
PoC relevantes de reclasificacion y reutilizacion de historial ya pasan sobre
\texttt{zcash\_client\_sqlite}.

\subsection{Por que esta fase es importante}

La Fase 2 transforma un problema de admision transitoria en:

\begin{quote}
\textbf{corrupcion persistente de estado wallet}
\end{quote}

Una vez persistido el estado, el sistema puede reconstruir objetos
internamente plausibles aunque el anchor final de verdad on-chain ya no sea el
que manda.

\section{Fase 3: corrupcion de intent en outputs enviados}

\subsection{Bug 3A: PCZT output intent fields son wallet-bound}

\paragraph{Hallazgo confirmado.}
\texttt{user\_address} y la metadata propietaria de output no estan protegidos
por consenso, sighash ni binding signatures.

pero el wallet los sigue tratando como si fuesen verdad fuerte luego de
extraer la transaccion.

\subsection{Bug 3B: full UA mutable despues de elegir receiver concreto}

\paragraph{Hallazgo confirmado.}
Una vez que proposal/build eligio un receiver concreto de una UA, el resto de
la UA ya no esta tx-bound; aun asi, la capa de extraction/storage puede
guardar y mostrar la UA completa desde metadata mutable.

\subsection{Bug 3C: sender recovery solo protege memo, no recipient/value/account}

\paragraph{Hallazgo confirmado.}
En ciertos paths de sent-output reconstruction desde PCZT, sender recovery se
usa para memo, mientras recipient, value y account salen primero de metadata
mutable.

\subsection{Consecuencia global de la fase 3}

Esta fase no cambia la tx on-chain, pero si cambia a quien cree el wallet que
le pagaste, que full UA cree que se uso, que cuenta interna cree que estuvo
involucrada y que valor cree que esta asociado a cierto output sent-history.

\section{Fase 4: amplificacion en spendability local}

\subsection{Como las fases anteriores contaminan gasto local}

Si el wallet ya persistio note internals no revalidados, nullifier metadata
caller-supplied, tree positions no rebindeadas y recipient o account semantics
falsas, entonces las capas posteriores pueden reportar spendable balance
equivocado, intentar generar witness sobre posiciones inconsistentes,
considerar notas como spendables hasta que falle una verificacion tardia y
generar propuestas o builds que solo fallan mas adelante.

\subsection{Tipo de fallo esperado}

No esperamos aqui un gasto on-chain valido con nullifier mal ligado. Lo que
esperamos es \textbf{late failure}, \textbf{user-visible inconsistency} y
\textbf{wallet-state corruption}.

\section{Fase 5: amplificacion por reorg}

\subsection{Bug 5A: deep rewind no limpia todo al target pedido}

\paragraph{Hallazgo confirmado.}
En deep rewind, el wallet puede preservar witness state, mined tx state y
nullifier-derived spentness por encima del target solicitado, hasta el pruning
floor.

\subsection{Por que esto agrava la cadena}

Si el estado ya venia contaminado por fases previas, el deep rewind puede
hacer que esa contaminacion dure mas tiempo, mezclar estado ``viejo pero
preservado'' con scan queue ``nueva pero aun no reparada'' y reforzar
decisiones locales de spendability sobre una base criptografica inconsistente.

\section{Fase 6: DoS como cierre de la cadena}

\subsection{Bug 6A: valores fuera de rango paniquean}

\paragraph{Hallazgo confirmado.}
Ciertos note values out-of-range pueden paniquear al convertir a
\texttt{Zatoshis}.

\subsection{Bug 6B: compact metadata malformada paniquea scanning}

\paragraph{Hallazgo confirmado.}
Txid, nullifiers y otros campos compact protobuf malformados pueden disparar
panic en scanning.

\subsection{Valor ofensivo de esta fase}

Si queremos una PoC fuerte de impacto local, el DoS es valioso porque puede
cortar el flujo de autoreparacion por rescan, dejar la base o el estado
intermedio en un punto donde el usuario solo ve que ``el wallet esta roto'' y
convertir corrupcion semantica en problema operacional visible.

\section{Ataques compuestos mas prometedores}

\subsection{Cadena A: input no validado + persistencia + spendability}

\begin{enumerate}[leftmargin=2em]
    \item alimentar al wallet una tx decryptable pero no plenamente validada;
    \item lograr que se persista;
    \item observar que metadata de nota / nullifier / posicion no queda
    suficientemente reatada;
    \item intentar un path que use ese estado como spendable;
    \item demostrar fallo tardio o inconsistencia local visible.
\end{enumerate}

\subsection{Cadena B: PCZT mutable + historial falso + persistencia}

\begin{enumerate}[leftmargin=2em]
    \item crear un PCZT valido;
    \item mutar metadata output-side;
    \item extraer y almacenar;
    \item probar que el wallet guarda recipient/value/account/history distinto
    del comprometido por la tx;
    \item intentar que esa asociacion falsa sobreviva reprocesamiento local.
\end{enumerate}

\subsection{Cadena C: estado contaminado + deep rewind}

\begin{enumerate}[leftmargin=2em]
    \item introducir estado wallet-side inconsistente;
    \item forzar deep rewind;
    \item observar que parte del estado criptografico sobrevive arriba del
    target;
    \item probar que spendability o history siguen reflejando ese estado mixto.
\end{enumerate}

\section{Impacto compuesto final}

La forma mas honesta de describir el impacto es:

\begin{quote}
No tenemos un exploit de theft on-chain. Tenemos una \textbf{cadena de
integridad y disponibilidad de wallet local} que puede producir:
\end{quote}

historial falso, recipient intent falso, semantica de output errada, resumenes
historicos de transaccion que heredan esa misma reclasificacion falsa,
balances o spendability local equivocadas, fallos tardios al gastar y crash
durante scanning, recovery o reprocesamiento.

Por eso la severidad correcta del compuesto es:

\textbf{High} para wallet local, \textbf{No Critical} y
\textbf{No direct funds theft confirmed}.

\section{Que partes queremos PoCear primero}

\subsection{Estado actual del target SQLite}

\texttt{zcash\_client\_sqlite} ya no esta bloqueado para las rutas principales
que importan en esta historia. Hoy ya confirma:

\begin{itemize}[leftmargin=2em]
    \item sent history falso por metadata mutable de PCZT;
    \item reclasificacion falsa de output externo como
    \texttt{InternalAccount};
    \item reutilizacion de esa reclasificacion en superficies de tx history;
    \item y la ruta legacy-context $\rightarrow$ sent history visible.
\end{itemize}

\subsection{Targets obligatorios}

Las PoC deben hacerse del mismo modo que la investigacion previa: como tests
en backend (\texttt{zcash\_client\_backend}), como tests en memory
(\texttt{zcash\_client\_memory}) y, cuando se trate de persistencia e
historial wallet-side concreto, tambien sobre
\texttt{zcash\_client\_sqlite}.

\subsection{Razon de esos targets}

\texttt{zcash\_client\_backend} nos deja demostrar el bug en la capa logica o
API. \texttt{zcash\_client\_memory} nos deja demostrar persistencia y
comportamiento wallet-side con iteracion rapida. Y
\texttt{zcash\_client\_sqlite} ya nos deja confirmar que la misma corrupcion
alcanza tambien al store relacional real y a sus superficies SQL de sent
history / tx history.

\section{Plan de PoC recomendado}

\subsection{PoC 1: version ``minima fuerte''}

\paragraph{Objetivo.}
Demostrar que una tx o estado shielded no plenamente validado puede ser
aceptado por APIs wallet-side bajo semantica de contexto equivocada y, desde
ahi, servir como base para un flujo compuesto de corrupcion de estado.

\paragraph{Forma.}
Hoy esta implementada como test en \texttt{zcash\_client\_backend} y ya esta
enganchada con una PoC wallet-visible en
\texttt{zcash\_client\_memory}.

\subsection{PoC 2: version ``PCZT + history corruption''}

\paragraph{Objetivo.}
Demostrar recipient, value, account e history corruption a partir de metadata
mutable PCZT en un flujo local end-to-end.

Esta PoC ya quedo implementada y validada en esta branch, y sigue siendo el
lugar mas natural para profundizar. La infraestructura conceptual ya existe;
backend y \texttt{zcash\_client\_memory} muestran la misma distincion entre
recipient comprometido y recipient persistido; y el ataque compuesto puede
apoyarse sobre esta misma base para sumar mas semantica errada, no solo
recipient.

\subsection{PoC 3: version ``reorg amplification''}

\paragraph{Objetivo.}
Demostrar que, una vez contaminado el estado local, un deep rewind no lo limpia
totalmente y deja spendability o witness state en condicion inconsistente.

\paragraph{Estado actual.}
Esta linea esta documentada y confirmada por auditoria, pero
\textbf{bloqueada en memory} por falta de implementacion de
\texttt{rewind\_to\_chain\_state}. Por lo tanto debe tratarse como siguiente
etapa compuesta, no como prerequisito inmediato.

\subsection{PoC 4: version ``panic as amplifier''}

\paragraph{Objetivo.}
Cerrar la historia mostrando que ciertas entradas malformadas pueden evitar
que el wallet se autorepare mediante rescan o recovery.

\section{Que no debemos exagerar en la PoC}

Para que la PoC sea solida, no conviene afirmar robo directo de fondos,
bypass de consenso ni doble interpretacion on-chain de una misma tx.

La PoC tiene que concentrarse en lo que si esta soportado por los hallazgos:
corrupcion de estado wallet local, persistencia de metadata o note state mal
reatada, spendability local equivocada, historial o intent falsos y DoS.

\section{Resumen final}

El viejo documento estaba bien para un solo bug PCZT. Ya no alcanza.

La imagen tecnica correcta hoy es la siguiente. \texttt{librustzcash} tiene
varios bugs wallet-side medianos o bajos, y varios de ellos comparten la misma
debilidad: confiar demasiado en datos no plenamente validados o no
cryptographically rebound. Al combinarlos, el impacto real sube. El mejor
target de demostracion es un \textbf{wallet local}, y la forma correcta de
PoCearlo en esta branch sigue siendo con tests en backend y memory, aceptando
que la parte de reorg hoy no puede cerrarse en memory sin implementar primero
soporte faltante del target.

Ese es el ataque sobre el que conviene construir la siguiente tanda de PoC
compuesta, arrancando por las rutas que ya estan demostradas y encadenandolas
antes de intentar cubrir la parte de reorg.

\end{document}
